\documentclass[11pt]{amsart}
\usepackage[T1]{fontenc}
\usepackage[utf8]{inputenc}
\usepackage{amsmath,amssymb,amsthm}
\usepackage{mathtools}
\usepackage[margin=1.1in]{geometry}

\providecommand{\C}{\mathbb{C}}
\providecommand{\P}{\mathbb{P}}
\providecommand{\R}{\mathbb{R}}
\providecommand{\Z}{\mathbb{Z}}
\providecommand{\dbar}{\bar{\partial}}
\providecommand{\g}{\mathfrak{g}}
\providecommand{\gBPS}{\mathfrak{g}_{\mathrm{BPS}}}
\providecommand{\Obs}{\mathrm{Obs}}
\providecommand{\Sym}{\mathrm{Sym}}
\providecommand{\CE}{\mathrm{CE}}
\providecommand{\chir}{\mathrm{chir}}
\providecommand{\Ran}{\mathrm{Ran}}
\providecommand{\Hom}{\mathrm{Hom}}
\providecommand{\cL}{\mathcal{L}}
\providecommand{\cM}{\mathcal{M}}
\providecommand{\cE}{\mathcal{E}}
\providecommand{\cD}{\mathcal{D}}
\providecommand{\cO}{\mathcal{O}}
\providecommand{\cT}{\mathcal{T}}
\providecommand{\cA}{\mathcal{A}}
\providecommand{\cF}{\mathcal{F}}
\providecommand{\CoHA}{\mathrm{CoHA}}
\providecommand{\Fact}{\mathrm{Fact}}
\providecommand{\kappach}{\kappa_{\mathrm{ch}}}
\providecommand{\kappaBKM}{\kappa_{\mathrm{BKM}}}
\providecommand{\kappacat}{\kappa_{\mathrm{cat}}}

\theoremstyle{plain}
\newtheorem{theorem}{Theorem}[section]
\newtheorem{proposition}[theorem]{Proposition}
\newtheorem{lemma}[theorem]{Lemma}
\newtheorem{conjecture}[theorem]{Conjecture}
\theoremstyle{definition}
\newtheorem{definition}[theorem]{Definition}
\newtheorem{construction}[theorem]{Construction}
\theoremstyle{remark}
\newtheorem{remark}[theorem]{Remark}

\title{Chiral CE of holomorphic Lie algebroids: five channels on
$\C^3$ and $K3 \times E$}
\author{Wave 13 G4/20}
\date{}

\begin{document}
\maketitle

\section{The target}
E4/20 constructed $\CE^*_{\dbar,\chir}(\cL,\cM)$, the chiral
Chevalley--Eilenberg complex on a complex threefold. We now apply this
functor to five specific holomorphic Lie algebroids, each recovering a
concrete factorization algebra and refining one face of the
CY-to-chiral correspondence.

\section{Cycle 1: Tangent Lie algebroid $\cT_{\C^3}$}
\begin{construction}[Tangent algebroid]\label{con:tangent}
Let $\cT_{\C^3} = \Omega^{0,\bullet}(\C^3, T^{1,0}\C^3)$ with
$\dbar$-differential and bracket given on coefficients by the
Schouten--Nijenhuis bracket of holomorphic vector fields; anchor
$\cT_{\C^3} \to \mathrm{Der}(\cO_{\C^3})$ is the natural action.
\end{construction}
\begin{proposition}\label{prop:tangent-factorization}
$\CE^*_{\dbar,\chir}(\cT_{\C^3}, \cO_{\C^3})$ is the factorization
algebra of holomorphic observables in the Costello--Li
$\beta\gamma$--$bc$ system with fields valued in $\cT_{\C^3}$; its
factorization homology on an open polydisc $D^3 \subset \C^3$ computes
$\Sym^\bullet\!\bigl(H^{0,\bullet}_c(D^3, T^{1,0})[-1]\bigr)^\vee$,
equivalently the BV pushforward of the infinitesimal holomorphic
diffeomorphism algebroid.
\end{proposition}
\textbf{Attack.} The Schouten bracket has no purely chiral restriction
to $\Ran^{\leq 2}\C^3$: it is a differential operator of order one on
both inputs, not a point-supported chiral bracket. \textbf{Heal.} In
the $\cD$-module formalism (\cite[\S3.5]{BD04}), a first-order bracket
is encoded by a chiral operation with support enlarged to a
first-infinitesimal neighbourhood of the diagonal; Francis--Gaitsgory
\cite{FrancisGaitsgory12} show this is equivalent to a chiral bracket
on the jet-Lie algebroid $J^\infty\cT_{\C^3}$, on which
$d_{\chir}^2 = 0$ holds exactly.

\section{Cycle 2: Current algebra $\cE = \Omega^{0,\bullet}(\C^3,\g)$}
\begin{proposition}[6D hCS observables]\label{prop:6d-hcs}
For finite-dimensional Lie algebra $\g$ and
$\cE = \Omega^{0,\bullet}(\C^3,\g)$ with pointwise bracket on
coefficients,
\[
  \Obs^{cl}_{\mathrm{hCS}}(\C^3)
   \;\simeq\; \CE^*_{\dbar,\chir}(\cE, \cO_{\C^3})
\]
as factorization algebras on $\C^3$, matching E4/20
Proposition~\ref{prop:hcs-obs-chiral-ce-recall} below and Costello's
6D hCS identification \cite{CostelloGaiotto18}.
\end{proposition}
\textbf{Attack.} Pointwise bracket violates the generic-point support
condition of BD chiral brackets. \textbf{Heal.} The
$\delta$-distribution along the diagonal is precisely the $\Delta_*$
image in Definition~\ref{def:chiral-ce-cochains-recall}; it is a
legitimate chiral bracket in the $\cD$-module sense because
$\Delta_*\cE \hookrightarrow j_*j^*(\cE \boxtimes \cE)$ in the canonical
short exact sequence.

\section{Cycle 3: Koszul duality at $d=3$}
\begin{proposition}[Chiral Koszul duality]\label{prop:koszul-3d}
Let $U^{\chir}_{\dbar}(\cE)$ denote the Dolbeault-enhanced chiral
envelope: the $E_1$-algebra on the holomorphic line times the
Dolbeault totalisation along transverse directions. Then
\[
  U^{\chir}_{\dbar}(\cE)
    \;\simeq\; \CE^*_{\dbar,\chir}(\cE, \C)^{!},
\]
with $(-)^{!}$ the Koszul dual in the $(\infty,1)$-category of
$E_3$-algebras \cite[Thm.~1.10]{Francis13}, specialised to the
holomorphic context of Kapranov--Vasserot
\cite[\S5]{KapranovVasserot14}.
\end{proposition}
\textbf{Attack.} Francis' Koszul duality requires $E_n$-algebra inputs;
$\cE$ a priori carries only an $E_1$ plus Dolbeault structure.
\textbf{Heal.} The Dolbeault structure upgrades $E_1$ to $E_3$ by the
framed-$E_n$ formalism of Ayala--Francis \cite{AyalaFrancis15}:
$\Omega^{0,\bullet}(\C^3, -)$ is a framed $\C^3$-factorization
algebra, automatically $E_3$ on the underlying topological site.
This is the $(\infty,1)$-categorical lane; the chain-level lane is the
explicit double complex of Definition~\ref{def:dolbeault-chiral-ce-recall}.

\section{Cycle 4: Factorization homology and CoHA}
\begin{conjecture}[CoHA--chiral CE]\label{conj:coha-chiral-ce}
For $\cE = \Omega^{0,\bullet}(\C^3,\g)$ (equivalently, for $X = \C^3$
viewed as CY$_3$),
\[
  \int_{\C^3} \CoHA(\C^3)
   \;\simeq\; \CE^*_{\chir}\!\bigl(U^+\!\bigl(\gBPS(\C^3)\bigr), \C\bigr),
\]
with $\CoHA(\C^3)$ the Kontsevich--Soibelman cohomological Hall algebra
of $\C^3$, $U^+$ its positive half, and $\gBPS(\C^3)$ the BPS Lie
algebra.
\end{conjecture}
\textbf{Attack.} $\CoHA(\C^3) = Y^+$ (positive half), not the full
Yangian (CLAUDE.md key facts). Does factorization homology of a
positive half make categorical sense? \textbf{Heal.} $Y^+$ is an
augmented $E_1$-algebra in the chiral Hall $\infty$-category of
Kapranov--Vasserot \cite[\S7]{KapranovVasserot14}. Augmentation
$Y^+ \to \C$ supplies the unit for
$\int_{\C^3}(-)$; on the other side, $\gBPS(\C^3)$ sits inside
$\CoHA$ as primitive elements (Davison--Meinhardt), identifying
$U^+(\gBPS) = \CoHA(\C^3)$ as Hopf algebras.
Conjecture~\ref{conj:coha-chiral-ce} then restates Francis--Ayala
Koszul duality in the $\chir$-enhanced setting; the obstacle is
morphism preservation of $\Phi: \mathrm{CY}\text{-cat}_3 \to
\mathrm{ChirAlg}$, which remains at scope-declaration status per
Pattern 273.

\section{Cycle 5: $K3 \times E$ and the chiral BKM}
\begin{construction}[$K3 \times E$ Lie algebroid]\label{con:k3xe-algebroid}
Let $X = K3 \times E$. The BPS Lie algebra $\gBPS(X)$ is conjecturally
the Monster / Borcherds BKM (\cite{Borcherds95}; compare
\texttt{cy\_d\_kappa\_stratification.tex},
Thm.~\texttt{thm:borcherds-weight-kappa-BKM-universal}). Set
$\cE_X := \Omega^{0,\bullet}(X, \gBPS(X))$, a holomorphic Lie
algebroid on the compact CY$_3$ $X$ with trivial anchor.
\end{construction}
\begin{conjecture}[Chiral BKM factorization algebra]\label{conj:chiral-bkm}
$\cF^{\chir}_{\mathrm{BKM}}(X)
  := \CE^*_{\dbar,\chir}(\cE_X, \cO_X)$
is the factorization algebra of the putative chiral BKM on
$K3 \times E$, with
\[
  \kappach\!\bigl(\cF^{\chir}_{\mathrm{BKM}}\bigr)
    = \chi\!\bigl(\CE^*_{\dbar,\chir}(\cE_X, \cO_X)\bigr)
    = 0,
\]
matching $\kappacat(K3\times E) = 0$ (total space, not fibre), whereas
the fibre-restricted complex $\CE^*_{\dbar,\chir}(\cE_X|_{K3},
\cO_{K3})$ contributes the Mukai-enhanced Heisenberg weight
$\kappaBKM(\Phi_1) = c_1(0)/2 = 5$.
\end{conjecture}
\textbf{Attack.} Is $\gBPS(K3\times E)$ well-defined as a holomorphic
Lie algebra, given that the six routes to $G(K3\times E)$ are six
different constructions (CLAUDE.md) and $G$ is unconstructed in
general? \textbf{Heal.} Conjecture~\ref{conj:chiral-bkm} is stated
with \texttt{ClaimStatusConjectured}: $\gBPS(X)$ is defined here as the
Davison--Meinhardt primitive piece of $\CoHA(X)$ \cite{DavisonMeinhardt20},
which exists as an $E_1$-Lie object without requiring the BKM
identification. The identification with Borcherds' monster algebra is
conjectural; the $\CE^*_{\dbar,\chir}$-construction of the
factorization algebra is not.
\textbf{Second attack.} Does $\kappach = 0$ clash with the spectrum
$\{2,3,5,24\}$ on $K3\times E$? \textbf{Heal.} No: $\{2,3,5,24\}$
come from four distinct constructions (Mukai lattice, Igusa
$\Phi_{10}$ via Gritsenko $\Delta_5$, BKM Borcherds weight, K3
fibre-rank), each tagged by its own subscript label. The total-space
$\kappach$ is $\chi(\cO_{K3\times E}) = \chi(\cO_{K3})\,\chi(\cO_E) =
2 \cdot 0 = 0$ by Künneth multiplicativity; the $\{2,3,5,24\}$ values
belong to $\kappacat(K3) = 2$, $\kappa_{\mathrm{fiber}} = 3$,
$\kappaBKM(\Phi_1) = 5$, and the K3 fibre-rank $24$ respectively.

\section{Recalls from E4/20}
\begin{definition}[E4 recall]\label{def:chiral-ce-cochains-recall}
$\CE^n_{\chir}(\cL, \cM) :=
 \Hom_{\cD_{X^n}}\!\bigl(\Sym^n_{\chir}\cL,\; \Delta^{(n)}_*\cM\bigr)$;
total differential $d_{\chir}$ via Jacobi on $\Ran^{\leq 3}$.
\end{definition}
\begin{definition}[E4 recall]\label{def:dolbeault-chiral-ce-recall}
$\CE^*_{\dbar,\chir}(\cL,\cM)
   = \bigl(\bigoplus_n \Omega^{0,\bullet}(X^n, \Hom(\Sym^n_{\chir}\cL,
     \cM))^{\Sigma_n},\, \dbar + d_{\chir}\bigr)$.
\end{definition}
\begin{proposition}[E4 recall]\label{prop:hcs-obs-chiral-ce-recall}
$\Obs^{cl}_{\mathrm{hCS}}(\C^3) \simeq
  \CE^*_{\dbar,\chir}(\cE, \cO_{\C^3})$ for
$\cE = \Omega^{0,\bullet}(\C^3,\g)$.
\end{proposition}

\begin{thebibliography}{9}
\bibitem{BD04} A.~Beilinson and V.~Drinfeld, \emph{Chiral Algebras},
  AMS Colloquium Publications 51, 2004, \S3.4--3.7.
\bibitem{Borcherds95} R.~Borcherds, \emph{Automorphic forms on
  $O_{s+2,2}(\R)$ and infinite products}, Invent.\ Math.\ 120 (1995),
  161--213.
\bibitem{KapranovVasserot14} M.~Kapranov and E.~Vasserot, \emph{The
  cohomological Hall algebra of a surface and factorization
  cohomology}, arXiv:1901.07641, \S5 and \S7.
\bibitem{Francis13} J.~Francis, \emph{The tangent complex and
  Hochschild cohomology of $E_n$-rings}, Compositio Math.\ 149 (2013),
  430--480.
\bibitem{FrancisGaitsgory12} J.~Francis and D.~Gaitsgory,
  \emph{Chiral Koszul duality}, Selecta Math.\ 18 (2012), 27--87.
\bibitem{AyalaFrancis15} D.~Ayala and J.~Francis, \emph{Factorization
  homology of topological manifolds}, J.~Topol.\ 8 (2015), 1045--1084.
\bibitem{CostelloGaiotto18} K.~Costello and D.~Gaiotto, \emph{Twisted
  supergravity and its quantization}, arXiv:1812.01110.
\bibitem{DavisonMeinhardt20} B.~Davison and S.~Meinhardt, \emph{Cohomological
  Donaldson--Thomas theory of a quiver with potential, and
  quantum enveloping algebras}, Invent.\ Math.\ 221 (2020), 777--871.
\bibitem{LurieHA} J.~Lurie, \emph{Higher Algebra}, 2017,
  \S5.5 ($E_n$-algebras) and \S7.3 (Koszul duality).
\end{thebibliography}

\end{document}
